\begin{corollary}[candidate nonvanishing for \texorpdfstring{$\operatorname{Pic}^0$}{Pic0} without topology]
\label{cor:candidate-pic0-route}
\Verified\ For the constructed threefold and every
$L\in\operatorname{Pic}^0(X)$,
\[
 H^0(W,\Omega_X^1|_W\otimes A_L)\ne0,
 \qquad
 H^2(W,(T_X\otimes L)|_W)\ne0,
\]
and
\[
 R^2f_*(T_X\otimes L)\ne0.
\tag{5.17f}
\]
This proof uses the analytic construction of the cusp fibre but no global
homology, Mayer--Vietoris, or Leray calculation.
\end{corollary}

\begin{proof}
The cusp theorem \cref{thm:candidate-cusp-geometry} supplies the full compact
reduced normal-crossings fibre, its $\mathrm{dP}_6$ normalization, and the
trivialization of $\omega_X$ near $W$.  By
\cref{prop:pic0-pullback}, $\eta^*A_L$ is trivial, so
\cref{prop:conductor} gives the first nonvanishing.  Apply
\cref{prop:cartier-duality-absolute} to $E=T_X\otimes L$; its dual space is
exactly the preceding nonzero section space, giving the fibrewise $H^2$
nonvanishing.  The fibre has complex dimension two, so its coherent
$H^3$ vanishes.  The global construction makes $X$ compact and $f$
nonconstant; \cref{lem:curve-flatness} therefore gives properness and
$f$-flatness of $T_X\otimes L$.  The base-change isomorphism of
\cref{prop:analytic-base-change} now identifies the fibre of the $R^2$
stalk with the nonzero $H^2$ group.  Hence the stalk, and therefore the
direct-image sheaf, is nonzero.
\end{proof}

\begin{remark}[reconciliation of the three restriction routes]
The Pic$^0$ argument proves (5.17f) without global homology,
Mayer--Vietoris, or Leray.  The all-Pic statement
\cref{prop:leray-edge-restriction} avoids citing the packaged theorem
$H^2(X;\mathbb Z)=0$, but it reuses the same verified global integral
Leray and specialization machinery from which that theorem is deduced.
By contrast, \cref{prop:direct-boundary-picard} supplies a second all-Pic
proof without any global Leray differential: it recovers the primitive
$4-3=1$ obstruction directly from the finite logarithmic gluing and the
cusp boundary.  It still uses the local integral cusp-specialization
theorem.  The third proof in \cref{thm:analytic-all-picard} replaces both
that specialization and the cusp Chern-class bounding-chain evaluation by
logarithmic $\mathcal O^*$ factors, the argument principle, and the exact
$\Psi$-equivariant support function.  Thus the purely analytic
strengthening is now complete; the three proofs have genuinely different
dependency sets.
\end{remark}

\subsection{Abstract conductor theorem for the candidate fibre}

